\documentclass[aspectratio=169]{beamer}

\usepackage[T1]{fontenc}
\usepackage{lmodern}
\usepackage{microtype}

\usepackage{minted}
\usemintedstyle{emacs}
\setminted{bgcolor=gray!10}
\usepackage{xcolor}

\title{Lightweight Word Embeddings (Word2Vec) and CPU BERT Demos}
\author{Francis Bond}
\date{}


\begin{document}

\begin{frame}
  \titlepage
\end{frame}

% ----------------------------
\begin{frame}{What is Word2Vec?}
\begin{itemize}
  \item \textbf{Goal:} learn \textbf{word embeddings} (dense vectors) from raw text.
  \item \textbf{Distributional idea:} words that occur in similar contexts have similar vectors.
  \item Two classic training objectives:
    \begin{itemize}
      \item \textbf{CBOW:} predict the target word from surrounding context words.
      \item \textbf{Skip-gram:} predict surrounding context words from the target word.
    \end{itemize}
  \item Result: geometry in vector space supports similarity, clustering, and some analogies.
\end{itemize}
\end{frame}

% ----------------------------
\begin{frame}{How Word2Vec is trained (intuition)}
\begin{itemize}
  \item Slide a window over a corpus; generate (context, target) pairs.
  \item Train a small neural model to make good predictions of words from contexts (or vice versa).
  \item Use tricks for efficiency:
    \begin{itemize}
      \item \textbf{Negative sampling} (common): discriminate true pairs from random “noise” pairs.
      \item \textbf{Subsampling} frequent words to reduce dominance of \textit{the, of, and, \dots}
    \end{itemize}
  \item The trained weights become word vectors; similarity is often cosine similarity.
\end{itemize}

\vspace{0.5em}
\small
\textit{Practical note:} for classroom laptops, \textbf{loading a small pre-trained model} is usually easier than training on a large corpus.
\end{frame}

% ----------------------------
\begin{frame}[fragile]{Gensim: CPU-friendly embeddings}
\begin{itemize}
  \item \textbf{Gensim} provides efficient implementations of Word2Vec and tools for using
        pre-trained vectors (KeyedVectors).
  \item Typical workflow:
    \begin{enumerate}
      \item install (\texttt{pip install gensim})
      \item load a small pre-trained model (or train a tiny one for a demo)
      \item query similarity / nearest neighbors / analogies
    \end{enumerate}
  \item Keep it lightweight: use a \textbf{small model} (tens of MB) and CPU.
\end{itemize}

\begin{minted}{bash}
pip install gensim
\end{minted}
\end{frame}

% ----------------------------
\begin{frame}[fragile]{Gensim: loading vectors + similarity}
\begin{itemize}
  \item You can load vectors saved in Gensim format, or in word2vec format.
  \item Once loaded, you can ask for nearest neighbors and cosine similarity.
\end{itemize}

\begin{minted}{python}
from gensim.models import KeyedVectors

# Example (Gensim native format)
# kv = KeyedVectors.load("vectors.kv", mmap="r")

# Example (word2vec text format)
# kv = KeyedVectors.load_word2vec_format("vectors.txt", binary=False)

# Then:
# kv.most_similar("university", topn=10)
# kv.similarity("cat", "dog")
\end{minted}

\small
\textit{Tip:} For large vectors, \texttt{mmap="r"} keeps RAM usage down.
\end{frame}

% ----------------------------
\begin{frame}[fragile]{Compositional / analogy arithmetic (classic demo)}
\begin{itemize}
  \item The famous analogy style query:
  \[
    \vec{king} - \vec{man} + \vec{woman} \approx \vec{queen}
  \]
  \item In Gensim, use \texttt{most\_similar} with positive/negative sets.
\end{itemize}

\begin{minted}{python}
# "king" - "man" + "woman"  -> "queen" (often)
kv.most_similar(positive=["king", "woman"],
                negative=["man"],
                topn=5)
\end{minted}

\small
You’ll often get \texttt{queen} near the top (depending on model + vocabulary).
\end{frame}

% ----------------------------
\begin{frame}[fragile]{Hyponymy / “is-a” vector trick (works... sometimes)}
\begin{itemize}
  \item People sometimes try a relation vector like:
  \[
    \vec{is\_a} \approx \vec{animal} - \vec{dog}
  \]
  then apply it to a new hyponym: \(\vec{queen} + \vec{is\_a}\).
  \item This is \textbf{not a principled guarantee}: hierarchies are not consistently linear in Word2Vec space.
\end{itemize}

\begin{minted}{python}
# Build a crude "is-a" direction from one example pair
is_a = kv["animal"] - kv["dog"]

# Apply it to another specific term
kv.most_similar(positive=["queen", is_a], topn=10)
\end{minted}

\small
Sometimes you’ll see broader categories (\texttt{monarch}, \texttt{royalty}, \texttt{person}...), sometimes you won’t.
That’s the teaching point.
\end{frame}

% ----------------------------
\begin{frame}{Why antonyms are hard for Word2Vec (and friends)}
\begin{itemize}
  \item Word2Vec learns from \textbf{shared contexts}.
  \item Antonyms often occur in \textbf{very similar contexts}:
    \begin{itemize}
      \item \textit{hot coffee} / \textit{cold coffee}
      \item \textit{high temperature} / \textit{low temperature}
    \end{itemize}
  \item So embeddings can place antonyms \textbf{close together} even though meanings oppose.
  \item Consequence: “nearest neighbors” is \textbf{not} the same as “synonyms”.
\end{itemize}

\vspace{0.5em}
\small
Useful classroom exercise: compare nearest neighbors for a word and ask students to label each neighbor as synonym / related / antonym / topical.
\end{frame}

% ----------------------------
\begin{frame}{From static embeddings to contextual models}
\begin{itemize}
  \item Word2Vec gives \textbf{one vector per word type}.
  \item But many words are polysemous:
    \begin{itemize}
      \item \textit{bank} (river vs.\ finance)
      \item \textit{chestnut} (tree/nut vs.\ “old chestnut” = stale joke/idea)
    \end{itemize}
  \item Contextual models (BERT-family) produce \textbf{token embeddings}:
    the vector for \textit{bank} depends on the sentence.
\end{itemize}
\end{frame}

% ----------------------------
\begin{frame}[fragile]{CPU BERT in class: small models + simple pipelines}
\begin{itemize}
  \item Hugging Face \texttt{transformers} runs fine on CPU for small demos.
  \item Use \textbf{distilled} models to keep it manageable:
    \begin{itemize}
      \item \texttt{distilbert-base-uncased} (English)
      \item multilingual options exist too, but are often heavier
    \end{itemize}
\end{itemize}

\begin{minted}{bash}
pip install transformers torch --index-url https://download.pytorch.org/whl/cpu
\end{minted}

\begin{minted}{python}
from transformers import pipeline
fill = pipeline("fill-mask", model="distilbert-base-uncased")
\end{minted}

\small
On CPU it’s slower than GPU, but fine for short sentences.

The first time we run it, it downloads the model
\begin{verbatim}
~/.cache/huggingface/
└── hub/
    └── models--distilbert-base-uncased/
\end{verbatim}

\end{frame}

% ----------------------------
\begin{frame}[fragile]{BERT “wow” moment: context sensitivity}
\begin{itemize}
  \item Same word, different context $\Rightarrow$ different predictions.
\end{itemize}

\begin{minted}{python}
# Same surface form, different sense-cues in context
s1 = "I sat on the bank of the river and watched the water."
s2 = "I went to the bank to open a new account."

# Fill-mask expects a [MASK] token; we mask a nearby word to probe context.
print(fill("I sat on the bank of the river and watched the [MASK].")[:3])
print(fill("I went to the bank to open a new [MASK].")[:3])
\end{minted}

\small
Students see that the \textbf{same surrounding topic} drives very different completions.
\end{frame}

% ----------------------------
\begin{frame}[fragile]{Negation: a simple demo (and a warning)}
\begin{itemize}
  \item Negation is famously tricky.
  \item A quick classroom probe: compare predictions with vs.\ without \textit{not}.
\end{itemize}

\begin{minted}{python}
print(fill("A robin is a [MASK].")[:5])
print(fill("A robin is not a [MASK].")[:5])
\end{minted}

\small
\textbf{Warning:} masked-LM objectives do not “reason” logically.
Sometimes the negated sentence still suggests plausible categories (or even repeats the positive behavior).
That unpredictability is itself a useful discussion point.
\end{frame}

% ----------------------------
\begin{frame}[fragile]{Metaphor / idiom strength: “double-edged \ldots”}
\begin{itemize}
  \item BERT often completes conventional metaphors / idioms well because it has seen them in varied contexts.
\end{itemize}

\begin{minted}{python}
print(fill("His words were a double-edged [MASK].")[:5])
\end{minted}

\small
Often \texttt{sword} appears near the top. This contrasts with Word2Vec arithmetic:
Word2Vec can do some analogies, but it does not represent compositional phrase meaning the same way.
\end{frame}

% ----------------------------
\begin{frame}[fragile]{Bert demo code is in one file}
\begin{itemize}
  \item Put everything in a single script, e.g. \texttt{bert\_demo.py}
\end{itemize}

\begin{minted}{bash}
$ python bert_demo.py
\end{minted}
\end{frame}

\end{document}
%%% Local Variables: 
%%% coding: utf-8
%%% mode: latex
%%% TeX-PDF-mode: t
%%% TeX-engine: xetex
%%% TeX-command-extra-options: "--shell-escape"
%%% End: 
